\problema{Longest Consecutive Sequence}{128}{src/01-grafos/128-longest-consecutive.cpp}{M}

Nos dan un arreglo de enteros sin ordenar. Queremos la longitud de la
secuencia de enteros \emph{consecutivos} más larga (por ejemplo $3,4,5,6$),
sin importar el orden en que aparezcan en el arreglo. El algoritmo debe
correr en $O(n)$.

\paragraph{La idea, con un ejemplo de la vida real.} Imagina un montón de
niños en el patio, cada uno con un número escrito en la playera, pero
revueltos y sin formarse. Los niños con números seguidos (por ejemplo
$3,4,5,6$) quieren tomarse de la mano y formar una fila. Nuestra tarea es
encontrar la fila más larga que se pueda armar.

El truco para no marearnos: en lugar de empezar a contar desde cualquier
niño, solo empezamos desde el que va \emph{al frente} de su fila, es decir,
el que no tiene a nadie con el número justo anterior. Si un niño tiene el
$4$ y existe alguien con el $3$, ese niño no es el frente: dejamos que el
que esté más adelante empiece a contar. Así cada fila se cuenta una sola
vez. Para preguntar rápido ``¿existe un niño con tal número?'' guardamos
todos los números en un \emph{conjunto} (\texttt{unordered\_set}): es como
una caja mágica en la que preguntar si algo está adentro es instantáneo, sin
tener que revisar uno por uno.

\begin{center}
\ttfamily
\begin{tabular}{c}
nums = \{100, 4, 200, 1, 3, 2\}\\[4pt]
100 $\to$ ¿existe 99? no $\to$ inicio, longitud 1\\
4\ \ $\to$ ¿existe 3?\ \ sí $\to$ no es inicio, se ignora\\
200 $\to$ ¿existe 199? no $\to$ inicio, longitud 1\\
1\ \ $\to$ ¿existe 0?\ \ no $\to$ inicio, cuenta 1,2,3,4 $\to$ longitud \textbf{4}\\
3, 2 $\to$ no son inicio, se ignoran\\
\end{tabular}
\end{center}

\noindent El máximo es \textbf{4}, por la secuencia $1,2,3,4$.

\paragraph{Cómo lo resuelve el código, paso a paso.}
\begin{enumerate}
  \item Metemos todos los números en la caja mágica
        (\texttt{unordered\_set}). Los repetidos se funden en uno solo, y
        eso está bien: un número repetido no alarga la fila.
  \item Recorremos los números uno por uno. Para cada número \texttt{num}
        preguntamos: ¿existe \texttt{num - 1} en la caja? Si \emph{sí}
        existe, este número no es el frente de su fila, así que lo
        saltamos.
  \item Si \texttt{num - 1} \emph{no} existe, entonces \texttt{num} es el
        frente. Empezamos a contar hacia arriba: ¿está \texttt{num + 1}?,
        ¿y \texttt{num + 2}? Seguimos mientras el siguiente exista,
        contando cuántos niños se unen a la fila.
  \item Guardamos la longitud de la fila más larga que hayamos visto. Esa
        es la respuesta.
\end{enumerate}
El código completo (abre el archivo con clic en el título):

\lstinputlisting{src/01-grafos/128-longest-consecutive.cpp}

\complejidad{$O(n)$}{$O(n)$}

\paragraph{¿Qué significa esa complejidad?} $O(n)$ quiere decir que el
trabajo total crece igual que la cantidad de números: si hay el doble de
números, tardamos más o menos el doble, nada peor. Quizá te preocupe que la
fila se cuente muchas veces, pero no pasa: como solo el niño del frente
empieza a contar, cada fila se recorre \emph{una sola vez} completa. Sumando
todo, cada número se toca a lo mucho dos veces (una al revisar si es frente y
otra al contarlo dentro de su fila). Si en cambio dejáramos que
\emph{cualquier} niño empezara a contar, una fila larga se recontaría desde
cada uno de sus miembros y el trabajo se dispararía a $O(n^2)$, mucho más
lento.

\paragraph{Consejos para la entrevista (Amazon).} El error más común es
ordenar el arreglo primero; eso funciona, pero tarda $O(n \log n)$ y no
cumple el $O(n)$ que suelen pedirte de forma explícita. Explica que el filtro
de ``solo el frente cuenta'' es justo lo que evita el trabajo repetido: sin
él, el algoritmo se vuelve $O(n^2)$ cuando hay una sola fila muy larga.
Menciona los casos raros: arreglo vacío (la respuesta es $0$), números
repetidos (no alargan la fila, por eso usamos un conjunto que los funde, no
un \texttt{multiset} ni un conteo de frecuencias) y números negativos
(funcionan igual, al algoritmo no le importa el signo).
